\chapter{Pendahuluan}

\section{Latar Belakang}

Sub-bab ini berisi uraian tentang latar belakang atau justifikasi ilmiah (dari \textit{paper} dan jurnal) dari permasalahan yang akan diteliti, alasan penelitian dan penelitian yang pernah dilakukan sebelumnya terkait fenomena tersebut . Latar belakang yang baik harus memperkenalkan masalah secara jelas, menunjukkan urgensi penelitian, dan menjelaskan mengapa pendekatan yang diusulkan relevan untuk menyelesaikan masalah tersebut.

\textcolor{red}{\textbf{Panduan Sitasi:} Dalam menyitasi sumber referensi, gunakan \texttt{\textbackslash cite\{key\}} agar tampilan referensi menjadi (author, year), seperti ini \cite{brauer2001}. Namun, jika ingin berbentuk narasi, gunakan \texttt{\textbackslash textcite\{key\}}, contoh: \textcite{cheney2001} menyatakan bahwa... Hapus panduan ini setelah tidak diperlukan.}

\textcolor{red}{\textbf{Contoh Sitasi Berbagai Jenis Referensi:} Paragraf berikut menunjukkan contoh sitasi untuk setiap jenis referensi yang didukung template ini. Periksa bagian Daftar Pustaka untuk memastikan format output sudah sesuai panduan. Hapus paragraf ini setelah tidak diperlukan.}

Pemodelan matematika dalam biologi populasi telah dibahas secara komprehensif dalam beberapa buku teks \cite{brauer2001, cheney2001, dai1989}. Pendekatan integral Henstock-Kurzweil juga telah dikembangkan melalui berbagai publikasi \cite{lee2000, lee1989}. Dalam bidang kimia, analisis spektrum senyawa organik telah diterjemahkan ke dalam bahasa Indonesia \cite{creswell1982}, sementara perangkat lunak komputasi kimia juga telah tersedia \cite{anonim1992}.

Beberapa penelitian dalam jurnal ilmiah telah membahas metode kromatografi \cite{davis1995}, model mekanika kuantum \cite{dewar1985}, dan metabolisme kulit \cite{finnen1987}. Penelitian terdahulu dalam bentuk skripsi juga telah membahas sistem linear dan aplikasinya \cite{husna2002}. Selain itu, laporan penelitian tentang sintesis antibiotik telah dilakukan \cite{jumina2001}.

Dalam konteks seminar ilmiah, aplikasi permainan dinamis pada sistem deskriptor telah dipresentasikan \cite{salmah2006}. Di bidang farmasi, beberapa senyawa turunan azaindol telah dipatenkan sebagai antibiotik \cite{wang2003}. Sumber daring juga dapat dijadikan referensi, seperti kajian tentang fungsi kelas Baire \cite{leung2000}.

\vspace{5mm}
\noindent\textbf{Contoh latar belakang penelitian untuk Magister Ilmu Komputer:} \\
\noindent\fbox{%
	\parbox{\textwidth}{%
		
		\hspace{1cm} Deteksi objek merupakan salah satu tugas fundamental dalam bidang \textit{computer vision} yang memiliki peran penting dalam berbagai aplikasi seperti kendaraan otonom, sistem pengawasan cerdas, dan robotika. Dalam beberapa tahun terakhir, arsitektur berbasis \textit{deep learning} seperti YOLO (\textit{You Only Look Once}) telah menunjukkan performa yang sangat baik dalam mendeteksi objek secara \textit{real-time} \cite{brauer2001}. Namun, model-model tersebut umumnya memiliki ukuran yang besar dan membutuhkan sumber daya komputasi yang tinggi, sehingga sulit untuk diimplementasikan pada perangkat \textit{edge} dengan keterbatasan memori dan daya komputasi seperti NVIDIA Jetson Nano atau Raspberry Pi.
		
		\hspace{1cm} Beberapa teknik optimasi model seperti \textit{pruning}, \textit{quantization}, dan \textit{knowledge distillation} telah diusulkan untuk mengurangi ukuran dan kompleksitas model tanpa mengorbankan akurasi secara signifikan. Meskipun demikian, penerapan teknik-teknik tersebut pada arsitektur YOLOv8 yang merupakan versi terbaru dari keluarga YOLO masih belum banyak dieksplorasi, khususnya dalam konteks \textit{deployment} pada perangkat \textit{edge}. Oleh karena itu, penelitian ini bertujuan untuk mengembangkan model YOLOv8 yang dioptimasi melalui kombinasi \textit{structured pruning} dan \textit{post-training quantization} agar dapat melakukan deteksi objek secara \textit{real-time} pada perangkat dengan sumber daya terbatas.
		
	}%
}

\vspace{5mm}
Latar belakang di atas memperkenalkan masalah deteksi objek pada perangkat \textit{edge}, menunjukkan urgensi penelitian terkait keterbatasan komputasi, dan mengusulkan pendekatan optimasi model sebagai solusi. Latar belakang ini memberikan dasar yang kuat bagi perumusan masalah dan tujuan penelitian, memastikan bahwa hasil penelitian memiliki relevansi dan signifikansi bagi bidang ilmu komputer.

\section{Rumusan Masalah}

Sub-bab ini berisi poin-poin yang menjelaskan masalah atau isu yang akan diteliti dalam suatu penelitian. Ini mencakup formulasi masalah yang jelas dan spesifik serta mempertimbangkan latar belakang dan deskripsi masalah yang terkait pada sub-bab sebelumnya. Perumusan masalah yang baik akan membantu menentukan fokus penelitian dan memberikan dasar untuk tujuan penelitian.

Jika rumusan masalah lebih dari 1, maka ditulis dengan format \textit{list}. Contoh sebagai berikut:

\begin{enumerate}
	\item Rumusan masalah 1.
	\item Rumusan masalah 2.
\end{enumerate}

\section{Batasan Masalah}

Batasan masalah adalah pembatasan yang menentukan lingkup dan ruang lingkup suatu penelitian. Ini membantu membatasi masalah dan isu yang akan diteliti, menentukan metode dan teknik penelitian, serta membatasi variabel yang akan diuji. Batasan penelitian juga mencakup keterbatasan dan hambatan yang mungkin timbul selama proses penelitian. Tujuan utama dari batasan masalah adalah membuat penelitian terfokus dan terarah, serta memastikan hasil yang diperoleh dapat digunakan untuk menjawab pertanyaan penelitian secara efektif dan akurat.

Batasan masalah dapat ditulis dengan format \textit{list}. Contohnya adalah sebagai berikut:

\begin{enumerate}
	\item Batasan 1.
	\item Batasan 2.
\end{enumerate}

\section{Tujuan Penelitian}

Tujuan penelitian adalah menentukan sasaran atau target yang ingin dicapai melalui proses penelitian. Tujuan penelitian bisa beragam sesuai dengan bidang ilmu yang dipelajari, topik penelitian, dan permasalahan yang akan dicari solusinya. Tujuan perlu disinkronkan dengan rumusan masalah. Setiap rumusan masalah boleh memiliki lebih dari 1 tujuan.

\noindent Catatan: Tujuan penelitian harus jelas, spesifik, terukur, dan dapat dicapai dalam jangka waktu yang ditentukan.

\vspace{5mm}
\noindent\textbf{Contoh tujuan penelitian untuk Magister Ilmu Komputer:} \\
\noindent\fbox{%
	\parbox{\textwidth}{%
		Berikut adalah contoh tujuan penelitian yang sesuai dengan rumusan masalah ``Bagaimana mengoptimasi arsitektur YOLOv8 agar dapat melakukan deteksi objek secara \textit{real-time} pada perangkat \textit{edge} tanpa penurunan akurasi yang signifikan?'':
		\begin{enumerate}
			\item Mengembangkan model YOLOv8 yang dioptimasi menggunakan teknik \textit{structured pruning} untuk mengurangi kompleksitas model.
			\item Menerapkan \textit{post-training quantization} (INT8) pada model yang telah di-\textit{pruning} untuk memperkecil ukuran model.
			\item Mengevaluasi performa model yang dioptimasi terhadap \textit{baseline} menggunakan metrik mAP (\textit{mean Average Precision}) dan FPS (\textit{Frames Per Second}) pada dataset COCO 2017.
			\item Melakukan \textit{deployment} dan pengujian inferensi \textit{real-time} pada perangkat NVIDIA Jetson Nano.
		\end{enumerate}
	}%
}

\section{Manfaat Penelitian}

Manfaat penelitian didefinisikan sebagai manfaat yang diperoleh apabila penelitian telah selesai dilakukan. Manfaat penelitian pada umumnya berbentuk daftar. Manfaat penelitian dapat berupa manfaat bagi dunia akademik dan/atau masyarakat.

\begin{enumerate}
	\item Manfaat akademik
	\begin{itemize}
		\item Manfaat akademik 1.
		\item Manfaat akademik 2.
	\end{itemize}
	\item Manfaat praktis
	\begin{itemize}
		\item Manfaat praktis 1.
		\item Manfaat praktis 2.
	\end{itemize}
\end{enumerate}

\section{Kontribusi Penelitian}

Uraikan kontribusi penelitian yang dilakukan dibandingkan dengan penelitian sebelumnya yang dijadikan sebagai basis penelitian. Kontribusi dapat berupa metode baru, peningkatan performa, pendekatan yang belum pernah dilakukan, atau kombinasi teknik yang menghasilkan solusi yang lebih baik.
